\documentclass{article}
\usepackage{tikz}
\pagestyle{empty}
\newsavebox\fig
\newenvironment{figura}{\begin{lrbox}{\fig}}{\end{lrbox}%
  \paperwidth=\dimexpr\wd\fig+12pt\relax \paperheight=\dimexpr\ht\fig+\dp\fig+12pt\relax
  \pagewidth=\paperwidth \pageheight=\paperheight \hoffset=-1in \voffset=-1in
  \shipout\vbox{\kern6pt\hbox{\kern6pt\usebox\fig}}}
\usepackage{fontspec}
% CM Bright, la letra de las figuras de las Feynman Lectures, en su versión
% OpenType (CMU Bright, de cm-unicode). Las OTF están en figuras/fuentes/ y
% render_figuras.py las pone en OPENTYPEFONTS. Los rótulos salen en mayúsculas
% (Letters=Uppercase); las variables van con \var, en otra familia, y quedan
% en minúscula. \textit también pasa a mayúsculas: no usarlo para variables.
\setmainfont{cmunbmr.otf}[ItalicFont=cmunbmo.otf, BoldFont=cmunbsr.otf,
  BoldItalicFont=cmunbso.otf, Letters=Uppercase, Scale=0.88]
\newfontfamily\varfont{cmunbmo.otf}[Scale=0.88]
\setmonofont{Menlo}[Scale=0.85]
\newcommand\var[1]{{\varfont #1}}
% CMU Bright no tiene ≥, ≤ ni ≠: salen del modo matemático
\catcode`≥=\active \protected\def≥{\ensuremath{\geq}}
\catcode`≤=\active \protected\def≤{\ensuremath{\leq}}
\catcode`≠=\active \protected\def≠{\ensuremath{\neq}}
\usetikzlibrary{arrows.meta,positioning,calc,decorations.pathreplacing,shapes.misc}
\definecolor{tinta}{HTML}{2B2F36}
\definecolor{suave}{HTML}{8A9099}
\definecolor{acento}{HTML}{C0392B} % rojo de la pizarra: watermark, rechazos
\definecolor{relleno}{HTML}{EEF1F5}
\definecolor{segundo}{HTML}{2E7D6B} % verde de la pizarra: segunda serie
\tikzset{
  every picture/.style={line width=0.7pt, color=tinta, draw=tinta, text=tinta, >={Stealth[length=2.2mm,width=1.6mm]}},
  caja/.style={draw=tinta, fill=relleno, rounded corners=1.5pt, minimum width=11mm, minimum height=7mm, font=\small},
  vida/.style={draw=suave, line width=0.5pt, dashed},
  msg/.style={->, draw=tinta},
  etiq/.style={font=\footnotesize, inner sep=1.5pt},
  cod/.style={font=\ttfamily\footnotesize, inner sep=1.5pt},
  acentuado/.style={draw=acento, text=acento},
}
% El vocabulario de piezas de las figuras, para que todas dibujen igual lo
% mismo. Lo carga estilo.tex. Regla: lo que aparece en dos figuras o más vive
% acá, y lo que aparece en una sola se define en esa figura. Cada forma está
% dibujada con su nombre en catalogo.tex (render_figuras.py --catalogo).
\usetikzlibrary{decorations.pathmorphing, shapes.multipart, fit, backgrounds, matrix}

\tikzset{
  % --- nodos -------------------------------------------------------------
  % una réplica, un servidor o un proceso; se agranda solo si el rótulo no entra
  nodo/.style={circle, draw=tinta, fill=relleno, minimum size=9mm, inner sep=1pt, font=\small},
  % el líder de Raft: doble círculo
  lider/.style={nodo, double=relleno, double distance=1.2pt},
  % un componente chico adentro de otra cosa (los nodos dentro del sistema)
  componente/.style={circle, draw=tinta, fill=white, minimum size=5mm, inner sep=0pt},
  % varias máquinas apiladas en pisos: app / stub, app / Raft, árbol / log;
  % \node[capas=2] {aplicación \nodepart{two} stub};
  capas/.style={rectangle split, rectangle split parts=#1, draw=tinta, fill=relleno,
    rounded corners=1.5pt, minimum width=22mm, font=\small, align=center},
  capas/.default=2,
  % un grupo de nodos encerrado (quórum, preference list), detrás de ellos:
  % \begin{scope}[on background layer]\node[grupo, fit=(a)(b)] {};\end{scope}
  grupo/.style={draw=segundo, fill=segundo!10, rounded corners=4mm, inner sep=2mm},
  % un paso numerado: \node[paso] {1};
  % el recuadro de una máquina, un cliente o un servidor que contiene cajas;
  % se dibuja con un rectangle o con fit=(a)(b)
  marco/.style={draw=tinta, rounded corners=4pt},
  paso/.style={circle, draw=tinta, fill=white, inner sep=0.8pt, minimum size=4mm, font=\scriptsize},
  % --- tiempo ------------------------------------------------------------
  % un evento sobre una línea de tiempo; el color es opcional: evento=acento
  evento/.style={circle, fill=#1, inner sep=0pt, minimum size=5pt},
  evento/.default=tinta,
  % el mensaje ondulado de Lamport
  onda/.style={->, decorate, decoration={snake, amplitude=0.5mm, segment length=3.5mm,
    pre length=1mm, post length=2.5mm}},
  % una transacción o un intervalo de TrueTime: el segmento con topes;
  % \draw[intervalo] (0,0) -- (3,0);  \draw[intervalo, segundo] ...
  intervalo/.style={{Bar[width=2.6mm]}-{Bar[width=2.6mm]}},
  % --- tiras y tablas ----------------------------------------------------
  % una celda de una tira: el log, los snapshots (ver \tira)
  % outer sep=0: si no, los bordes de dos celdas vecinas quedan corridos y se ven dobles
  celda/.style={draw=tinta, fill=white, minimum height=9mm, inner sep=2pt, outer sep=0pt,
    font=\small, align=center, anchor=south west},
  % una tabla de la pizarra, sin bordes: \matrix[tabla] (t) {K & V\\ \var{x} & 10\\};
  % el filete bajo el encabezado va con \filete{t}{columnas}
  tabla/.style={matrix of nodes, nodes={font=\small, minimum height=6.5mm, inner sep=2pt},
    column sep=5mm},
  % --- flechas -----------------------------------------------------------
  % control y datos, como en las figuras de GFS: fina y gruesa
  control/.style={->, line width=0.5pt},
  datos/.style={->, line width=1.6pt},
  % la respuesta o el camino alternativo
  respuesta/.style={msg, dashed},
  % "de esto se pasa a esto", entre dos dibujos
  implica/.style={-{Implies}, double, double distance=2pt},
  % --- marcas ------------------------------------------------------------
  % el corte de un anillo, una clave sobre él, una marca gruesa sobre una línea
  marca/.style={line width=1.4pt},
  % la llave de una anotación; el argumento es la amplitud
  llave/.style={decorate, decoration={brace, amplitude=#1}, draw=suave},
  llave/.default=6pt,
  % --- piezas armadas: \pic at (x,y) {nombre}; ---------------------------
  % un rechazo, un mensaje perdido: la cruz roja
  % {rechazo=segundo} cambia el color, para una segunda falla
  pics/rechazo/.style={code={\draw[draw=#1, line width=1.4pt] (-0.16,-0.16) -- (0.16,0.16) (-0.16,0.16) -- (0.16,-0.16);}},
  pics/rechazo/.default=acento,
  % aceptado, al día: la marca verde, que arranca en el punto dado
  ok/.pic={\draw[segundo, line width=1.4pt] (0,0) -- ++(0.15,-0.17) -- ++(0.32,0.42);},
  % la partición de red: el rayo rojo, centrado en el punto dado
  particion/.pic={\draw[acento, line width=1pt] (0.15,0.75) -- (-0.08,0.1) -- (0.08,0.05) -- (-0.15,-0.75);},
  % una base de datos, un disco, un ledger: el cilindro con la base en el punto dado
  base/.pic={\def\w{0.5}\def\h{1.1}\def\e{0.15}
    \draw[fill=relleno] (-\w,\h) -- (-\w,0) arc[start angle=180, end angle=360, x radius=\w, y radius=\e] -- (\w,\h);
    \draw[fill=relleno!50!white] (0,\h) ellipse[x radius=\w, y radius=\e];},
  % una persona, el usuario: con los pies en el punto dado; \pic (p) deja (p-cabeza)
  persona/.pic={\draw (0,1.05) circle (0.2);
    \draw (0,0.85) -- (0,0.4) (-0.35,0.7) -- (0.35,0.7) (-0.28,0) -- (0,0.4) -- (0.28,0);
    \coordinate (-cabeza) at (0,1.05);},
}

% Tacha un nodo ya dibujado: el líder caído, la máquina que falló.
\newcommand\tachar[1]{\draw[acento, line width=1.4pt]
  (#1.north west) -- (#1.south east) (#1.north east) -- (#1.south west);}

% La flecha dentro de un texto, dibujada, del color del texto.
\newcommand\flecha{\colorlet{colorflecha}{.}\tikz[baseline=-0.55ex]\draw[->, draw=colorflecha] (0,0) -- (0.5,0);}

% El anillo de consistent hashing, centrado en el origen, con el corte en 0
% arriba; deja el radio en \R. Sus puntos van en sentido horario.
%   \anillo[radio]   \nodoanillo[color]{ángulo}{rótulo}{anchor}   \claveanillo{ángulo}
\newcommand\anillo[1][2.4]{\def\R{#1}\draw (0,0) circle (\R);
  \draw[marca] (90:\R-0.2) -- (90:\R+0.2);}
\newcommand\nodoanillo[4][tinta]{\fill[#1] (#2:\R) circle (2.4pt);
  \node[etiq, text=#1, anchor=#4] at (#2:\R+0.12) {#3};}
\newcommand\claveanillo[1]{\draw[segundo, marca] (#1:\R-0.22) -- (#1:\R+0.22);}

% Una muesca sobre un intervalo o una línea de tiempo: \muesca[segundo]{(2,0)}
\newcommand\muesca[2][]{\draw[#1] #2 ++(0,-1.3mm) -- ++(0,2.6mm);}

% Una operación de una transacción: \op{W}{x} es W con x de subíndice.
\newcommand\op[2]{\var{#1}\textsubscript{\var{#2}}}

% Historias de transacciones en el tiempo (clase 11). El eje x es el tiempo
% real; una fila por transacción, la primera arriba, cada 1,2 de alto:
%   \tx[x]{y}{1}                          el nombre T₁ a la izquierda, en x = 0
%   \transaccion[nodo]{y}{x0}{x1}{\op{W}{x} = 1}
%                                         el intervalo, con la operación arriba
%   \ts[below]{(x,y)}{10}                 el timestamp, en verde, en el instante
%                                         en que se asigna: commit o start
%   \intervalott[below]{y}{x0}{x1}{9}{\elegido{12}}
%                                         el intervalo de TrueTime pedido en ese
%                                         instante, con el valor elegido en caja
%   \draw[pedido] (x,y) -- (x,y');        la flecha que baja a buscarlo
%   \tiempo[tiempo real]{x0}{x1}{y}       el eje, en gris, al pie
%   \hora[below]{(x,y)}{10}               un tiempo real, en gris, donde ocurre
% Verde es lo que el sistema conoce (timestamps, TrueTime); gris, el tiempo
% real, que nadie conoce y que solo dibujamos para razonar.
\tikzset{
  pedido/.style={->, dashed, draw=segundo, line width=0.5pt},
  ts/.style={text=segundo, inner sep=1pt},
}
\newcommand\tx[3][0]{\node[anchor=west] at (#1,#2) {\var{T}\textsubscript{#3}};}
\newcommand\transaccion[5][op]{\draw[intervalo] (#3,#2) -- (#4,#2) node[midway, above] (#1) {#5};}
\newcommand\ts[3][below]{\draw[marca, draw=segundo] #2 ++(0,-1.5mm) -- ++(0,3mm);
  \node[ts, #1=5pt] at #2 {@\,#3};}
\newcommand\elegido[1]{\tikz[baseline=(e.base)]\node[draw=segundo, inner sep=1.5pt] (e) {#1};}
\newcommand\intervalott[6][below]{\draw[intervalo, draw=segundo] (#3,#2) -- (#4,#2);
  \node[ts, #1=4pt] at (#3,#2) {#5};
  \node[ts, #1=4pt] at (#4,#2) {#6};}
\newcommand\tiempo[4][tiempo real]{\draw[->, draw=suave, line width=0.5pt] (#2,#4) -- (#3,#4)
  node[etiq, text=suave, below left, yshift=-3pt] {#1};}
\newcommand\hora[3][below]{\draw[draw=suave, line width=0.6pt] #2 ++(0,-1.3mm) -- ++(0,2.6mm);
  \node[etiq, text=suave, #1=3pt] at #2 {#3};}

% Una tira de celdas, como el log: \tira[ancho]{nombre}{(origen)}{a, b, {x\\y}}.
% El origen es la esquina de abajo a la izquierda; deja los nodos nombre-1, …
\newcommand\tira[4][1.4]{\foreach \c [count=\i] in {#4}
  \node[celda, minimum width=#1cm] (#2-\i) at ($#3+({(\i-1)*#1},0)$) {\c};}

% El filete bajo el encabezado de una \matrix[tabla]: \filete{t}{columnas}
\newcommand\filete[2]{\draw ($(#1-1-1.south west)+(-2mm,0)$) -- ($(#1-1-#2.south east)+(2mm,0)$);}

% Reservadas, se diseñan con la primera clase que las use:
%   árbol (Zookeeper, Merkle) — clase 3

